\documentclass[11pt]{scrartcl}
\usepackage[sexy]{evan} %BLBLBLBLBLBLBLBLBLBLBLB EVANCHEN
\usepackage[T1]{fontenc}
\usepackage[utf8]{inputenc}
\usepackage{graphicx}
\usepackage{amsmath}
\usepackage{tikz}
\usepackage{tikz-cd}
\usepackage{asymptote}
\usepackage{amsthm}
\usepackage{amssymb}
\usepackage{gensymb}
\usepackage{amsfonts}
\usepackage[english,italian]{babel}
\usepackage{quoting}
\quotingsetup{font=big}
\usepackage{booktabs}
\usepackage{caption}
\usepackage{lmodern}
\usepackage{faktor}
\usepackage{bbm} %oer fare gli insiemi indicatori
\usepackage{leftindex} %per gli ideali e gli indici a sinistra
\usepackage{hyperref}
\usepackage{epigraph} 
\usepackage{pgfplots}
\usepackage[export]{adjustbox}

%Pacchetti per i codici
\usepackage{listings}
\usepackage{xcolor}
\usepackage{algpseudocode}

%Analisi e probabilità
%----------------------------------
\newcommand{\dif}{\text{d}}
%%% the pure symbol, see https://tex.stackexchange.com/a/718194/4427
\newcommand{\cind}{\perp\mspace{-9mu}\perp}
%%% define \indep
\NewDocumentCommand{\indep}{e{_}}{%
	\mathrel{%
		\IfNoValueTF{#1}{% no subscript
			\cind
		}{% subscript
			\mathchoice{\mathop{\cind}\limits_{#1}}{\cind_{#1}}{\cind_{#1}}{\cind_{#1}}%
		}%
	}%
}
%-----------------------------------

%Algebra lineare e geometria
%----------------------------------
\newcommand{\tr}[1]{\text{tr}\!\left(#1\right)}
\renewcommand{\rank}[1]{\text{rank}\!\left(#1\right)}
\newcommand{\Immagine}[1]{\text{Im}\left(#1\right)}
\renewcommand{\ker}[1]{\text{Ker}\left(#1\right)}
\renewcommand{\AA}{\mathbb{A}}
\newcommand{\LL}{\mathbb{L}}
\newcommand{\SO}{\text{SO}}
%----------------------------------

%Fisica
%----------------------------------
\newcommand{\UDiMisura}[1]{\,\text{#1}}
\newcommand{\ihat}{\hat{\imath}}
\newcommand{\jhat}{\hat{\jmath}}
\newcommand{\khat}{\hat{k}}
\newcommand{\posit}{\vec{r}}
\newcommand{\veloc}{\dot{\vec{r}}}
\newcommand{\accel}{\ddot{\vec{r}}}
%----------------------------------

%Informatica
%----------------------------------
\definecolor{codegreen}{rgb}{0,0.6,0}
\definecolor{codegray}{rgb}{0.5,0.5,0.5}
\definecolor{codepurple}{rgb}{0.58,0,0.82}
\definecolor{backcolour}{rgb}{0.95,0.95,0.92}

\lstdefinestyle{overleafwiki}{
	backgroundcolor=\color{backcolour},   
	commentstyle=\color{codegreen},
	keywordstyle=\color{magenta},
	numberstyle=\tiny\color{codegray},
	stringstyle=\color{codepurple},
	basicstyle=\ttfamily\footnotesize,
	breakatwhitespace=false,         
	breaklines=true,                 
	captionpos=b,                    
	keepspaces=true,                 
	numbers=left,                    
	numbersep=5pt,                  
	showspaces=false,                
	showstringspaces=false,
	showtabs=false,                  
	tabsize=2
}

\lstset{style=overleafwiki}
\newcommand{\bigo}[1]{\mathcal{O}\!\left(#1\right)}
\newcommand{\bigomega}[1]{\Omega\!\left(#1\right)}
\newcommand{\bigtheta}[1]{\Theta\!\left(#1\right)}
%----------------------------------

%Insiemistica e Algebra
%----------------------------------
\newcommand{\eset}{\varnothing}
\newcommand{\parti}[1]{\mathcal{P}\!\left(#1\right)}
\newcommand{\inj}{\hookrightarrow}
\newcommand{\surj}{\twoheadrightarrow}
\newcommand{\bij}{\hookrightarrow\mathrel{\mspace{-15mu}}\rightarrow}
\newcommand{\ordine}[2]{\text{ord}_{#1}\left(#2\right)}
\newcommand{\normal}{\mathrel{\unlhd}}
\newcommand{\MCD}{\text{MCD}}
\renewcommand{\mcm}{\text{mcm}}
%----------------------------------


\begin{document}
	\title{Appunti di Analisi 2}
	\subtitle{
		Davide Vittone (Modulo A) - UNIPD \\
		Annalisa Massaccesi (Modulo B) - UNIPD
	}
	\date{Settembre 2026 - Giugno 2027}
	
	\maketitle
	
	\tableofcontents
	
	\newpage
	
	\part{Modulo A}
	
	\section{Spazi metrici}
	
	\begin{definition}[Spazio Metrico]
		Uno spazio metrico è un insieme $X$ equipaggiato con una funzione $d:X\times X\to \RR$ tale che, per ogni $x,y,z\in X$ si abbia
		\begin{itemize}
			\item La distanza è sempre positiva, a meno che i punti coincidano, ovvero
			$$d(x,y)\ge 0$$
			con uguaglianza se e solo se $x=y$.
			\item La distanza è simmetrica, cioè
			$$d(x,y)=d(y,x).$$
			\item La disuguaglianza triangolare
			$$d(x,y)\le d(x,z)+d(z,y).$$
		\end{itemize}
		La funzione $d$ si chiama appunto \textit{distanza} dello spazio.
	\end{definition}
	Alcuni esempi sono
	\begin{example}[$\RR$ con varie distanze]
		L'insieme $\RR$ equipaggiato con 
		$$d(x,y):=\abs{x-y}$$
		è uno spazio metrico. Così come
		$$d(x,y):=\sqrt{\abs{x-y}}$$
		e in generale qualunque distanza della forma
		$$d(x,y):=\abs{x-y}^\alpha$$
		con $\alpha\in (0,1]$.
	\end{example}
	\begin{example}[$\CC$]
		L'insieme $\CC$ con la funzione
		$$d\left(z,w\right):=\abs{z-w}$$
		è uno spazio metrico.
	\end{example}
	Un esempio più strano è per esempio:
	\begin{example}[Distanza discreta]
		Qualunque insieme $X$ equipaggiato con la distanza
		$$d(x,y)=1-\delta_{xy}$$
		è uno spazio metrico.
	\end{example}
	Mentre uno ancora più classico è
	\begin{example}
		Lo spazio $\RR^n$ è uno spazio metrico con la distanza
		
	\end{example}
	
	\subsection{Spazi Normati e metrici}
	
	\begin{definition}[Spazio normato]
		Uno \textit{spazio normato} è uno spazio vettoriale su $\RR$ equipaggiato con una funzione \textit{norma} $\abs{\abs{\cdot}}:V\to [0,\infty)$ tale che, per ogni $x,y\in V$ e $\lambda\in \RR$ abbiamo
		\begin{itemize}
			\item La norma è sempre positiva, cioè
			$$\abs{\abs{x}}\ge 0$$
			con uguaglianza se e solo se $x$ è nullo.
			\item La norma è omogenea, ovvero
			$$\abs{\abs{\lambda x}}=\abs{\lambda}\abs{\abs{x}}.$$
			\item La norma è subadditiva, ovvero
			$$\abs{\abs{x+y}}\le \abs{\abs{x}}+\abs{\abs{y}}.$$ 
		\end{itemize}
		L'ultima proprietà in particolare può essere vista come un caso particolare della disuguaglianza triangolare.
	\end{definition}
	
	\begin{proposition}[Ogni spazio normato è metrico]
		Se $V$ con la norma $\abs{\abs{\cdot}}$ è uno spazio normato, allora lo spazio $V$ equipaggiato con la distanza
		$$d(x,y):=\abs{\abs{x-y}}$$
		è uno spazi metrico.
	\end{proposition}
	\begin{proof}
		Il fatto che la distanza sia semidefinita positiva segue immediatamente dal fatto che la norma sia semidefinita positiva. La simmetria segue dalla omogeneità per scalari negativi. La parte interessante è la disuguaglianza triangolare, per cui abbiamo
		$$d(x,y)=\abs{\abs{x-y}}=\abs{\abs{\left(x-z\right)+\left(z-y\right)}}\le \abs{\abs{x-z}}+\abs{\abs{z-y}}$$
		usando la subadditività, da cui abbiamo la tesi.
	\end{proof}
	
	\begin{example}[$\ell^1$]
		Possiamo equipaggiare $\RR^2$ con la norma
		$$\abs{\abs{(x_1,x_2)}}_1=\abs{x_1}+\abs{x_2}$$
		per renderlo uno spazio normato.
	\end{example}
	\begin{example}[$\ell^\infty$]
		Possiamo equipaggiare $\RR^2$ con la norma
		$$\abs{\abs{(x_1,x_2)}}_\infty=\max\left\{x_1,x_2\right\}$$
		per renderlo uno spazio normato.
	\end{example}
	Abbiamo visto a geometria che tutti gli spazi con prodotto interno sono anche normati, e di conseguenza sono anche metrici. A questo proposito ricordiamo la disuguaglianza di Cacuchy-Schwarz:
	\begin{theorem}[Cauchy-Schwarz]
		In tutti gli spazi con prodotto $V$ abbiamo, per ogni $x,y$:
		$$\abs{\langle x,y\rangle}\le \abs{\abs{x}}\abs{\abs{y}}$$
		con uguaglianza se e solo se $x,y$ sono linearmente dipendenti.
	\end{theorem}
	\begin{proof}
		Una dimostrazione alternativa consiste nel considerare il polinomio:
		$$P(t):=\abs{\abs{x+ty}}^2=\abs{\abs{x}}^2+2t\langle x,y\rangle +t^2\abs{\abs{y}}$$
		ma essendo che questo polinomio è sempre maggiore o uguale di $0$ abbiamo che il suo determinante deve essere non positivo, ovvero
		$$0\ge \Delta = 4\langle x,y\rangle^2-4\abs{\abs{x}}^2\abs{\abs{y}}^2$$
		che era la tesi.
	\end{proof}
	Da questa è immediata dimostrare che $\RR^n$ è uno spazio normato.
	\begin{corollary}[$\RR^n$ è metrico]
		Lo spazio $\RR^n$ equipaggiato con la distanza indotta dalla norma indotta dal prodotto scalare standard è metrico.
	\end{corollary}
	\begin{proof}
		Basta mostrare che è normato, le prime due proprietà seguono dalle proprietà del prodotto scalare, la subadditività invece segue da Cauchy-Schwarz:
		$$\abs{\abs{x+y}}^2=\abs{\abs{x}}^2+2\langle x,y\rangle +\abs{\abs{y}}^2\le \abs{\abs{x}}^2+2\abs{\abs{x}}\abs{\abs{y}}+\abs{\abs{y}}^2=\left(\abs{\abs{x}}+\abs{\abs{y}}\right)^2$$
		da cui la subadditività, e quindi la tesi.
	\end{proof}
	In uno spazio metrico ha senso definire:
	\begin{definition}[Palla aperta]
		In uno spazio metrico $(X,d)$ per ogni $x\in X$ e $r\in\RR$ non negativo definiamo la palla aperta di raggio $r$ come
		$$B(x,r):=B_r(x):=\left\{y\in X:d(x,y)<r\right\}.$$
	\end{definition}
	\begin{example}[Ferrovie francesi]
		Possiamo equipaggiare $\RR^2$ con la seguente distanza per renderlo uno spazio metrico:
		$$d(x,y)=
		\begin{cases}
			\abs{\abs{x-y}} & \text{se $x$, $y$ sono linearmente dipendenti} \\
			\abs{\abs{x}}+\abs{\abs{y}} & \text{altrimenti.}
		\end{cases}
		$$
		Che forma hanno le palle in questo spazio metrico? Beh se $r<\abs{\abs{x}}$ abbiamo che $B_r(x)$ è semplicemente un segmento centrato in $x$ di ampiezza $2r$ e allineato con $0$. D'altra parte se $r\ge \abs{\abs{x}}$ si aggiunge una palla centrata nell'origine di raggio $r-\abs{\abs{x}}$.
	\end{example}
	Possiamo definire un concetto di \textit{sottostruttura} in uno spazio metrico, ovvero
	\begin{definition}[Spazio metrico restrizione]
		Uno spazio metrico $(X,d)$ può essere \textit{ristretto} ad uno spazio $Y\subseteq X$ restringendo la funzione $d_{\mid Y\times Y}$. Le palle in questo caso soddisfano
		$$B_Y(x,r)=B_X(x,r)\cap Y.$$
	\end{definition}
	
	\subsubsection{Limiti}
	
	Uno spazio metrico ci dà tutti gli strumenti per parlare degli oggetti che abbiamo incontrato finora in analisi, in particolare 
	\begin{definition}[Successione]
		Una successione in uno spazio metrici $(X,d)$ è una funzione $x_n:\NN\to X$, denotata con $\left(x_n\right)_{n\in\NN}$. 
	\end{definition}
	\begin{definition}[Limite]
		Una successione $a_n$ a valori in $X$ si dice convergere a un punto $x\in X$ se vale
		$$\lim_{n\to\infty}d(x_n,x)=0$$
		abbiamo quindi tradotto la nozione di limite su spazi arbitrari mutuandola da $\RR$, su cui sapevamo lavorare. In questo caso allora possiamo scrivere $x_n\to x$ oppure
		$$\lim_{n\to\infty}x_n = x.$$
	\end{definition}
	\begin{proposition}[Unicità del limite]
		In qualunque spazio metrico $(X,d)$ se una successione $x_n$ ha limite $x$, questo è unico.
	\end{proposition}
	\begin{proof}
		Supponiamo che la successione abbia due limiti $x$ e $y$, allora 
		$$0\le d(x,y)\le d(x,x_n)+d(x_n,y)$$
		per ogni $n\in\NN$. D'altra parte facendo andare $n\to\infty$ abbiamo che entrambi gli addendi al RHS vanno a $0$, dunque 
		$$0\le d(x,y)\le 0$$
		da cui $x=y$ per fedeltà della distanza.
	\end{proof}
	Siamo interessati a lavorare soprattutto in $\RR^n$, ovvero:
	\begin{proposition}[Limiti in $\RR^m$]
		Sia $X=\RR^m$ con la distanza euclidea, definiamo la successione
		$$\mathbf{x}_n=\left(x_n^1,x_n^2,\dots,x_n^m\right)$$
		allora 
		$$\lim_{n\to\infty}x_n={\mathbf{x}}=\left({x}^1,{x}^2,\dots,{x}^n\right)$$
		se e solo se
		$$\lim_{n\to\infty}x_n^i=x^i$$
		per ogni $i$.
	\end{proposition}
	\begin{proof}
		Un verso è facile, abbiamo
		$$\abs{x_n^i-x^i}\le \abs{\abs{x_n-x}}=\sqrt{\sum_{i=1}^{m}\left(x_n^i-x^i\right)^2}$$
		da cui abbiamo che se la successione converge in $\RR^n$ ogni componente converge individualmente a quello che deve convergere. D'altra parte il contrario è molto facile da dimostrare in modo analogo, e quindi abbiamo finito.
	\end{proof}
	Come fatto su $\RR$ abbiamo:
	\begin{definition}[Punto di accumulazione]
		Sia $(X,d)$ è uno spazio metrico e $A\subseteq X$. Un punto $x_0\in X$ si dice \textit{punto di accumulazione} di $A$ se $\forall r>0$ si ha
		$$\left[A\cap B(x_0,r)\right]\setminus\left\{x_0\right\}\ne \eset.$$
	\end{definition}
	\'E facile mostrare ora, forse assumendo l'assioma della scelta purtroppo
	\begin{proposition}[Approssimabilità dei punti di accumulazione]
		Se $x_0$ è un punto di accumulazione di $A$ allora esisterà una successione $x_n\in A$ tale che $x_n\ne x_0$ per ogni $n$ e tale che $x_n\to x_0$.
	\end{proposition}
	Da cui possiamo finalmente definire la nozione di limite di funzione:
	\begin{definition}[Limite di funzione]
		Dati due spazi metrici $(X,d_X)$ e $(Y,d_Y)$ e una funzione 
		$$f:A\subseteq X\to Y$$
		diremo che, dato un punto di accumulazione $x_0$ di $A$, il limite di $f$ per $x$ che tende a $x_0$ si scrive come
		$$\lim_{x\to x_0}f(x)=y_0\in Y$$
		se e solo se, per ogni $\varepsilon>0$ esiste un $\delta>0$ tale che tutti i punti $\delta$-vicini a $x_0$ vengono mandati da $f$ in punti $\varepsilon$-vicini a $y_0$.
	\end{definition}
	E possiamo anche ritrovare il teorema-ponte dato ad Analisi 1:
	\begin{theorem}[Caratterizzazione sequenziale del Limite]
		Se $X,Y,f,A,x_0$ sono insiemi e funzioni come sopra allora sono equivalenti le seguenti:
		\begin{description}
			\item[A)] Per $x\to x_0$ la funzione $f$ tende a $y_0$, ovvero
			$$\lim_{x\to x_0}f(x)=y_0.$$ 
			\item[B)] Per ogni successione $x_n\to x$ in $X$ si ha che la successione $f(x_n)$ in $Y$ tende a $y_0$.
		\end{description}
	\end{theorem}
	\begin{proof}
		Dimostriamo che A implica B. Dato un $\varepsilon>0$ possiamo trovare un $\delta>0$ tale che 
		$$d(x_n,x_0)< \delta \implies d(f(x_n),y_0)<\varepsilon$$
		d'altra parte presa una $x_n\to x_0$ abbiamo che definitivamente si avrà $d(x_n,x_0)<\delta$, allora definitivamente si avrà
		$$d(f(x_n),y_0)<\varepsilon$$
		che è proprio la definizione di $f(x_n)\to y_0$. \\
		Dimostriamo ora il contrario, procedendo per assurdo. Se per assurdo esistesse un $\overline{\varepsilon}>0$ tale che non esista alcun $\delta>0$ per cui la condizione del limite si verifica, per definizione di limite allora abbiamo che
		$$d(x_n,x_0)<\frac{1}{n}$$
		è definitivamente vera, d'altra parte
		$$d(f(x_n),x_0)\ge \overline{\varepsilon}$$
		è frequentemente vera, il che è assurdo, e dunque abbiamo finito.
	\end{proof}
	
	\newpage
	
	\part{Modulo B}
	
\end{document}
